% Source: https://artofproblemsolving.com/community/c3h1115327_2015_national_mathcounts_27_and_29_and_target_4
% Stable identity: aops-topic-1115327-q04
Place $D=(0,0)$, $C=(12,0)$, $B=(12,12)$, $A=(0,12)$, so $M=(6,0)$.

\textbf{The crease.} Folding $A$ onto $M$ means the crease is the perpendicular bisector of $AM$. The midpoint of $AM$ is $(3,6)$ and $AM$ has slope $\dfrac{0-12}{6-0}=-2$, so the crease has slope $\dfrac{1}{2}$:
\[
y-6=\dfrac{1}{2}(x-3),\qquad\text{i.e.}\qquad y=\dfrac {x}{2}+\dfrac{9}{2} .
\]
It meets $AD$ ($x=0$) at $P=(0,4.5)$ and $BC$ ($x=12$) at $Q=(12,10.5)$. (Check: $PA=7.5=PM$, since $PM^2=6^2+4.5^2=56.25$.)

\textbf{The trapezoid.} The folded part $ABQP$ has the parallel vertical sides $AP=12-4.5=7.5$ and $BQ=12-10.5=1.5$, at distance $AB=12$ apart. Its area is
\[
\dfrac{7.5+1.5}{2}\cdot12=\boxed{54}.
\]
After folding it becomes the congruent trapezoid $PQB'M$ (shaded), where $B'$ is the image of $B$:

\begin{asy}
usepackage("amsmath");
size(6cm);
defaultpen(fontsize(10pt));
pair A=(0,12),B=(12,12),C=(12,0),D=(0,0),M=(6,0),P=(0,4.5),Q=(12,10.5);
pair Bp=reflect(P,Q)*B;
fill(P--Q--Bp--M--cycle,paleyellow);
draw(A--B--C--D--cycle,gray+dashed);
draw(D--C--Q--P--cycle,linewidth(0.9));
draw(P--Q--Bp--M--cycle,linewidth(0.9));
draw(A--M,red+dotted);
dot(M);dot(P);dot(Q);
label("$A$",A,NW);label("$B$",B,NE);label("$C$",C,SE);label("$D$",D,SW);
label("$M$",M,S);label("$P$",P,W);label("$Q$",Q,NW);label("$B'$",Bp,E);
label("$6$",(3,0),S);label("$6$",(9,0),S);label("$4.5$",(0,2.25),W);label("$7.5$",(0,8.25),W);label("$1.5$",(12,11.25),E);
label("\mbox{}",(min(currentpicture,user=true)+max(currentpicture,user=true))/2,black); // keeps the text after the figure black when LaTeX typesets the labels inline
\end{asy}

(The unfolded part $PQCD$ is the other trapezoid, of area $144-54=90$.)
